\documentclass[11pt]{article}
\usepackage{amsmath,amssymb,amsthm,geometry,hyperref}
\geometry{margin=1in}
\newtheorem{lemma}{Lemma}
\newtheorem{proposition}{Proposition}
\newtheorem{theorem}{Theorem}
\newtheorem{corollary}{Corollary}
\newcommand{\OO}{\mathcal O}
\newcommand{\ZZ}{\mathbb Z}
\newcommand{\QQ}{\mathbb Q}
\newcommand{\ff}[2]{#1^{\underline{#2}}}
\begin{document}
\title{Five-adic block transfer and index-interpolation obstruction for the ordinary quintic}
\author{Internal research note}
\date{17 November 2026}
\maketitle

For the ordinary support $X^2(A+BX^3)$, this note retains an exact
all-index Frobenius block transfer and a different, analytic
index-family calculation on the circle $v_5(u)=1$.
The block transfer preserves seven physical shape offsets and
constant-unit seeds, but does not yet preserve the required
joint three-cyclic rank relation.
The cancelled analytic family has exactly two common five-adic
indices. It differs from the actual finite moment polynomials
by a small, generally nonzero, resonant tail.
A hypothetical ordinary collision on that circle must consequently
be exponentially close to one of those two analytic indices.
An explicit constant-unit seed also gives a rank-one generalized
mixed pair at $u=10$, refuting a seed-blind strengthening of the
missing descent. Its seed differs from the actual seed $h=1$.
None of these results proves all-terminal quintic SFC or identifies an
ordinary integer counterexample.

\section{Exact Artin--Hasse transfer for seven physical offsets}
Let $K/\mathbb Q_5$ be a complete valued extension, let $\OO$ be its
valuation ring, and normalize $v(5)=1$.  Work on $v(z)>0$ with
\[
 \theta_z=z\frac{d}{dz},\qquad
 A=F(z)-\theta_z,\qquad
 F(z)=z+\sum_{r\ge1}a_rz^{5^r},\quad a_r\in5\OO .
\]
The transfer is evaluated in the inner region $v(z)>1/5$, so that
$v(w)>0$ after the substitution below.
A seed means $h\in\OO[[z]]$ whose reduction modulo the maximal ideal
is a nonzero constant.  For integers $n\ge0$ and $\delta$, define the commuting polynomial in $A$
\[
 P_{n,\delta}(A)=
 \prod_{k=0}^{n-1}(A+5n+\delta-1-3k),\qquad P_{0,\delta}=1 .
\]
Define $B_{0,\gamma,3}=1$ and
$B_{r+1,\gamma,3}=(z+3r+\gamma+2)B_{r,\gamma,3}-\theta_zB_{r,\gamma,3}$.
For $F=z$ and $h=1$, the expression $P_{n,\delta}(A)h$ specializes exactly to
$B_{n,\,2n+\delta,\,3}$, since reversing the factors gives
$\prod_{k=0}^{n-1}(z-\theta_z+2n+\delta+2+3k)1$.

Put
\[
 E(z)=\exp\left(\sum_{r\ge0}\frac{z^{5^r}}{5^r}\right),\qquad
 b=a_1-1,\qquad w=\frac{bz^5}{5}.
\]
The Artin--Hasse series $E$ and $E^{-1}$ have coefficients in
$\mathbb Z_5$.  For completeness, the coefficient identity
\[
 E(z)=\prod_{\substack{d\ge1\\5\nmid d}}
                 (1-z^d)^{-\mu(d)/d}
\]
follows by taking formal logarithms: the coefficient of $z^N$ in
the logarithm is $N^{-1}\sum_{d\mid N,\,5\nmid d}\mu(d)$,
which is $1/N$ if $N$ is a power of five and zero otherwise.
All exponents belong to $\mathbb Z_5$, so the factors and the
locally finite product are integral.  Consequently the gauge is
an analytic unit on the open unit disk.  Its first five inverse
coefficients are $[z^j]E^{-1}=(-1)^j/j!$ for $0\le j<5$.

\begin{lemma}[Exact gauge and potential closure]\label{lem:quintic-ah-gauge}
The gauge gives
\[
 E^{-1}AE=-\theta_z+\sum_{r\ge1}(a_r-1)z^{5^r}.
\]
On the residue-$j$ subspace $z^j\OO[[z^5]]$, after $w=bz^5/5$,
this operator is $-j+5L$, where
\[
 L=-\theta_w+w+\sum_{r\ge1}a'_rw^{5^r},\qquad
 a'_r=\frac{(a_{r+1}-1)5^{\,5^r-1}}{b^{5^r}}\in5\OO .
\]
In fact $v(a'_r)\ge 5^r-1$.  Thus the same class of potentials
is preserved.
\end{lemma}
\begin{proof}
The derivative of $\log E$ is
$\theta_z\log E=\sum_{r\ge0}z^{5^r}$.
The residue-$j$ action of $-\theta_z$ is
$-j-5\theta_{z^5}$.  Substitute $z^5=5w/b$.
The linear term is $(a_1-1)z^5/5=w$, and the remaining
coefficients are exactly those displayed.  Since $b\equiv-1\pmod5$,
it is a unit.
\end{proof}

\begin{proposition}[Seven-offset transfer]\label{prop:quintic-seven-offset}
Write $n=5\ell+s$, $0\le s<5$.  For $0\le j<5$ set
\[
 \begin{split}
 k_j&=\langle2(\delta-1-j)\rangle_5,\qquad
 \beta_j=\frac{j+3k_j-\delta+1}{5},\\
 \epsilon_j&=\mathbf1_{k_j<s},\qquad t_j=\ell+\epsilon_j,\\
 \delta'_j&=s-\beta_j+1-5\epsilon_j .
 \end{split}                                                    \tag{F1}
\]
Here $\langle\,\rangle_5$ is the least residue in $\{0,\ldots,4\}$.
Decompose the gauged seed exactly as
\[
 E^{-1}h=\sum_{j=0}^4z^jH_j(z^5).
\]
Define
\[
 \begin{split}
 c_{j,k}&=5n+\delta-1-j-3k,\\
 U_j(L)&=\prod_{\substack{0\le k<n\\k\not\equiv k_j\ (5)}}
                   (c_{j,k}+5L),\\
 h_j(w)&=U_j(L)H_j(5w/b).
 \end{split}
\]
Then every $h_j$ is again a seed and there is the exact analytic
identity
\[
 E^{-1}P_{n,\delta}(A)h
   =\sum_{j=0}^4 z^j5^{t_j}
                     P_{t_j,\delta'_j}(L)h_j(w).                 \tag{F2}
\]
Every $\delta\in\{-4,-3,-2,-1,0,1,2\}$ has all its descendants
$\delta'_j$ in the same set.  For $n\ge2$ every nonzero descendant
index $t_j$ is strictly smaller than $n$.
\end{proposition}
\begin{proof}
By Lemma~\ref{lem:quintic-ah-gauge}, the block action is $-j+5L$.
The factors divisible by five in block $j$ are exactly
$k=k_j+5a$, $0\le a<t_j$.  Each such factor is
\[
 -j+5L+5n+\delta-1-3(k_j+5a)
       =5(L+n-\beta_j-3a).
\]
The definition of $\delta'_j$ gives
$5t_j+\delta'_j-1=n-\beta_j$.  Hence their product is
$5^{t_j}P_{t_j,\delta'_j}(L)$.  The other factors are exactly
$U_j(L)$.  These are all polynomials in the same $L$, so they
commute.  This proves (F2) before evaluation on any cyclic fiber.

The constant term of $H_j$ reduces to
$h(0)(-1)^j/j!$, a unit.  The substitution $5w/b$ makes all
higher terms vanish modulo the maximal ideal.  For each
nonresonant factor $c_{j,k}$ is a five-adic unit and
$c_{j,k}+5L\equiv c_{j,k}$ modulo five as an operator on
$\OO[[w]]$.  Thus $h_j$ still has constant-unit reduction.

For closure assume $-4\le\delta\le2$.  Since $\beta_j$ is an
integer and
$-1\le j+3k_j-\delta+1\le21$, one has $0\le\beta_j\le4$.
If $\epsilon_j=0$, then $s\le k_j$ and
\[
 -3\le\delta'_j=s+1-\beta_j
 \le k_j+1-\left\lceil(3k_j-1)/5\right\rceil\le2 .
\]
The final inequality is checked at the five integers $0\le k_j<5$.
If $\epsilon_j=1$, then $1\le s\le4$, $k_j\le s-1$, and
$\beta_j\le\lfloor(3s+6)/5\rfloor$.  Therefore
\[
 -4\le s-4-\lfloor(3s+6)/5\rfloor
        \le\delta'_j=s-4-\beta_j\le0 .
\]
This proves closure.  Finally
$t_j\le\lceil n/5\rceil<n$ for $n\ge2$.
\end{proof}

The smaller proposed set $\{-3,-2,-1,0,1,2\}$ is \emph{not}
closed.  Already $\delta=-3$, $s=1$, $j=1$ gives
$k_j=0$, $\beta_j=1$, $\epsilon_j=1$, hence $\delta'_j=-4$.
The seven-offset correction is indispensable.

For the original physical offset $\delta=1$, the descendants,
in column order $j=0,\ldots,4$, are
\[
\begin{array}{c|rrrrr}
s&\delta'_0&\delta'_1&\delta'_2&\delta'_3&\delta'_4\\\hline
0&1&-1&0&-2&-1\\
1&-3&0&1&-1&0\\
2&-2&1&-3&0&1\\
3&-1&2&-2&1&-3\\
4&0&-2&-1&2&-2
\end{array}.
\]
For the original potential $F=z$ one has $b=-1$ and
$w=-z^5/5$.  Thus the cyclic argument $u=z^3$ transforms to
\[
 u'=w^3=-u^5/125,\qquad v(u')=5v(u)-3.                            \tag{F3}
\]
At later stages $u'=b^3u^5/125$, with the same valuation map.

\paragraph{Exact scope and remaining obstruction.}
This is an all-index operator and seed-closure identity, not a
joint-rank descent.  The five descendant seeds depend on $j,n,\delta$
through different polynomials $U_j(L)$; retaining only their
constant-unit reductions loses their exact coupling.
Moreover a rank relation at the three points
$z,\omega z,\omega^2z$ does not separately annihilate the five
residue blocks in (F2).  Since $5$ is coprime to $3$,
$w(\omega z)=\omega^2w(z)$ permutes the same three cyclic points,
but the five prefactors $z^j5^{t_j}$ remain coupled.
No implication from a rank drop of the original ordinary mixed
pair to a rank drop of a smaller physical mixed pair is proved.
The normalized inner region $v(u)>3/5$ therefore remains open.
No squarefreeness-preservation, SFC, independent/external review,
formal verification, or novelty claim is made.


\section{The cancelled analytic family and its exact scope}
Normalize $v_5(5)=1$.  The actual reverse ordinary moment polynomial is
\[
 \widehat M_n(u)=
 \sum_{j=0}^{n}\binom nj\frac{(5n-3j)!}{(5n)!}u^j,\qquad n\ge1.       \tag{I1}
\]
Indeed it is a nonzero scalar times
$u^n T_{n,\,2n+1,\,3}(u^{-1})$.
Thus consecutive ordinary quintic moments have a common reciprocal
root $u$ precisely when $\widehat M_m(u)=\widehat M_{m+1}(u)=0$.
We consider only the circle $v_5(u)=1$, so write $u=5w$ with $w$ a unit.

For every $j\ge1$ put
\[
 t_j=\left\lceil\frac{3j}{5}\right\rceil,\qquad
 b_j(M)=
 \frac{5^{j-t_j}\ff{(M-t_j)}{j-t_j}}
 {j!\displaystyle\prod_{\substack{1\le r<3j\\5\nmid r}}(5M-r)},
 \qquad b_0=1,                                                     \tag{I2}
\]
and define the \emph{cancelled analytic family}
\[
                  Q(M,w)=\sum_{j\ge0}b_j(M)w^j.                    \tag{I3}
\]
It is essential that (I3) is an infinite analytic continuation, not
the finite polynomial (I1).

\begin{lemma}[Analyticity and exact finite-polynomial comparison]\label{lem:quintic-cancelled-tate}
Each $b_j$ belongs to $\ZZ_5\langle M\rangle$, with coefficientwise bound
\[
 v_5(b_j)\ge e_j:=j-t_j-v_5(j!)\ge\frac{3j-11}{20}.                 \tag{I4}
\]
The $e_j$ are nonnegative and tend to infinity, so
$Q\in\ZZ_5\langle M,w\rangle$.
For an integer $m\ge1$, its first $m+1$ coefficients are exactly those
of $\widehat M_m(5w)$.  Moreover
\[
 b_j(m)=0\quad\text{if}\quad m<j\le\lfloor5m/3\rfloor.              \tag{I5}
\]
For every $m\ge8$ and every integral $w$,
\[
 v_5\bigl(Q(m,w)-\widehat M_m(5w)\bigr)\ge L_m,\qquad
 L_m=\max\{3,\lfloor m/4\rfloor\}.                                \tag{I6}
\]
\end{lemma}
\begin{proof}
Every factor $5M-r$ in (I2) has unit constant term and an inverse
in $\ZZ_5\langle M\rangle$.  Remove the $5$-part of $j!$;
the remaining prefactor has valuation $e_j$.
Since $t_j\le(3j+4)/5$ and
$v_5(j!)\le(j-1)/4$, (I4) follows.
Directly $e_1=e_2=0$, $e_3=1$, and (I4) makes every $e_j$
nonnegative for $j\ge4$.

For $j\le m$, factor the denominator of
\[
 \binom mj\frac{(5m-3j)!}{(5m)!}5^j
       =\frac{\ff m j\,5^j}{j!\,\ff{(5m)}{3j}} .
\]
Its factors divisible by five are exactly $5(m-a)$,
$0\le a<t_j$, and cancel the corresponding numerator factors.
This gives exactly (I2).  When $m<j$ and $t_j\le m$,
the remaining falling factorial in (I2) contains zero, proving (I5).
After (I5), every possible error index has $j>5m/3$.
By (I4), $e_j>m/4-11/20$, which implies
$e_j\ge\lfloor m/4\rfloor$.

It remains to check the lower bound three in (I6).
For $j=5q+s$, $0\le s<5$, the exact formula is
\[
 e_j=q-v_5(q!)+d_s,\qquad (d_0,d_1,d_2,d_3,d_4)=(0,0,0,1,1).
\]
For $j\ge13$ and $q=2$ the two possibilities $s=3,4$ give three.
For $q\ge3$, the bound
$q-v_5(q!)\ge(3q+1)/4$ gives at least three.
When $m\ge8$, the first possible error index is at least fourteen.
This proves (I6).
\end{proof}

\begin{lemma}[Exact first residuals]\label{lem:quintic-index-residuals}
The following are coefficientwise identities in the indicated
restricted power-series rings:
\[
 \begin{split}
 Q(M,w)&\equiv1-2w+2w^2\pmod5,\\
 Q(M,2)&\equiv25(M-1)^2\pmod{125},\\
 Q(M,9)&\equiv25(2M^2-M-2)\pmod{125}.
 \end{split}\tag{I7}
\]
In particular $Q(M,2)$ and $Q(M,9)$ are divisible by $25$.
For $c=2,9$ define
\[
 F_c(M,Y)=\frac{Q(M,c+25Y)}{25},\qquad
 G_c(M,Y)=\frac{Q(M+1,c+25Y)-Q(M,c+25Y)}{25}.                       \tag{I8}
\]
These belong to $\ZZ_5\langle M,Y\rangle$.  Their reductions are
\[
 \begin{array}{c|cc}
 c&\overline F_c&\overline G_c\\\hline
 2&(M-1)^2+Y&2M-1\\
 9&2M^2-M-2-Y&1-M .
 \end{array}                                                      \tag{I9}
\]
\end{lemma}
\begin{proof}
Lemma~\ref{lem:quintic-cancelled-tate} gives $e_j\ge3$ for $j\ge13$.
Expand the unit denominators in (I2) and use only $0\le j\le12$.
The following complete table gives $b_j(M)c^j$ modulo $125$:
\[
\begin{array}{c|l|l}
 j&c=2&c=9\\\hline
0&1&1\\
1&-50M^2-55M+1&25M^2-60M-58\\
2&-52&-53\\
3&-50M^2+40M-5&-25M^2+20M+60\\
4&25M^2-45M+35&25M^2+55M-15\\
5&35M^2+5M+45&-5M^2+35M-60\\
6&-50M^3-40M^2+35M-50&50M^3-10M^2+40M+50\\
7&-45M^2-5M+25&-35M^2+10M-50\\
8&25M^3+50M^2+50M&25M^3+50M^2+50M\\
9&50M^3-50M^2+50M-50&-25M^3+25M^2-25M+25\\
10&50M^4-50&50-50M^4\\
11&50M^4+50M^3+50M^2+50M&25M^4+25M^3+25M^2+25M\\
12&25M^4+50M^3-25M^2-50M&25M^4+50M^3-25M^2-50M
\end{array}
\]
Each entry follows directly from
$(5M-r)^{-1}\equiv-\sum_{a=0}^2(5M)^a/r^{a+1}\pmod{125}$.
Summing the columns gives (I7) coefficientwise, not merely at
integral test points. The retained symbolic verifier checks all
twenty-six literal entries and both shifted differences in $M$.
The first equation in (I7) has simple roots $w=2,4$ modulo five.
The other two equations imply the divisibility in (I8);
substitution of $c+25Y$ adds
$Y Q_w(M,c)$ modulo five after division.
At $c=2$ this derivative is $1$, and at $c=9$ it is $-1$.
Taking the difference between $M+1$ and $M$ gives (I9).
\end{proof}

\begin{theorem}[A precise index-interpolation obstruction]\label{thm:quintic-five-adic-index-obstruction}
There are exactly two pairs
\[
 (\alpha_2,\omega_2),\quad(\alpha_9,\omega_9)
              \in\ZZ_5\times\ZZ_5^\times
\]
with
\[
                    Q(\alpha,w)=Q(\alpha+1,w)=0.                  \tag{I10}
\]
They are characterized uniquely by
\[
 \begin{split}
 (\alpha_2,(\omega_2-2)/25)&\equiv(3,1)\pmod5,\\
 (\alpha_9,(\omega_9-9)/25)&\equiv(1,4)\pmod5 .
 \end{split}                                                      \tag{I11}
\]
Equivalently $(\alpha_2,\omega_2)\equiv(3,27)$ modulo $(5,125)$
and $(\alpha_9,\omega_9)\equiv(1,109)$ modulo $(5,125)$.
Thus a blanket noncollision theorem for all integral five-adic
indices of the natural family (I3) is false.
\end{theorem}
\begin{proof}
Apply Lemma~\ref{lem:quintic-index-residuals}.
Every unit root of $Q(M,w)$ reduces to $2$ or $4$ modulo five,
and its simple root lifts modulo twenty-five to $2$ or $9$,
respectively, by (I7).  Hence any solution of (I10) belongs to
one of the coordinates (I8).
In the first row of (I9), the only common zero over the residue
field is $(M,Y)=(3,1)$; in the second row it is $(1,4)$.
The Jacobian matrices with columns $(M,Y)$ at those points are
\[
 \begin{pmatrix}4&1\\2&0\end{pmatrix},
 \qquad
 \begin{pmatrix}3&-1\\-1&0\end{pmatrix}\quad\bmod5 ,
\]
with determinants $-2$ and $-1$.
At each successive precision, the two correction digits are
therefore determined by an invertible linear system over $\mathbb F_5$.
Completeness gives one lift in each residue disk and the same
linear argument gives uniqueness.  There are no other residual
solutions.
\end{proof}

\begin{corollary}[All-index necessary exponential closeness on this circle]\label{cor:quintic-integer-index-closeness}
Let $m\ge8$ be an ordinary integer.  Suppose
$\widehat M_m(u)=\widehat M_{m+1}(u)=0$ at $v_5(u)=1$,
and put $w=u/5$.  Then for exactly one branch $c=2$ or $9$,
with the notation above,
\[
 \begin{split}
 v_5(m-\alpha_c)&\ge L_m-2,\\
 v_5(w-\omega_c)&\ge L_m,\\
 v_5(u-5\omega_c)&\ge L_m+1 .
 \end{split}                                                       \tag{I12}
\]
In particular the only possible residue classes are
$m\equiv3\pmod5$ on the branch $w\equiv27\pmod{125}$,
and $m\equiv1\pmod5$ on the branch $w\equiv109\pmod{125}$.
\end{corollary}
\begin{proof}
Lemma~\ref{lem:quintic-cancelled-tate}, also at $m+1$, gives
$v_5(Q(m,w)),v_5(Q(m+1,w))\ge L_m$.
The residual calculation (I7) first places $w$ in
$c+25\OO$ for one of the two branches.
Set $Y=(w-c)/25$.
Then both entries of $(F_c(m,Y),G_c(m,Y))$ have
valuation at least $L_m-2\ge1$, so (I9) places $(m,Y)$
in the residue disk of its unique exact root from
Theorem~\ref{thm:quintic-five-adic-index-obstruction}.

For an integral restricted analytic map with Jacobian invertible
modulo five on one residue disk, Taylor expansion gives
\[
 \min_i v_5(\Phi_i(x)-\Phi_i(x_*))
       =\min_i v_5(x_i-x_{*,i}).
\]
Indeed the linear term is multiplication by an invertible
integral matrix, which preserves the minimum valuation, and
every higher term has strictly larger valuation because
$x-x_*$ has positive valuation.
Apply this to $\Phi=(F_c,G_c)$ and its exact root.
It gives both coordinate bounds at least $L_m-2$.
Rescaling $w=c+25Y$ and $u=5w$ proves (I12).
This argument remains valid in any complete algebraic extension
containing the putative root.
\end{proof}

\paragraph{Scope boundary.}
The family $Q$ is not the actual finite moment polynomial for
all integral or ordinary indices: its coefficients can resume
after the zero interval (I5).  The tail is small at any fixed
precision when $m$ is sufficiently large, but is not identically
zero.  Consequently (I10) does not identify an ordinary integer
counterexample, or even a common zero of two actual finite
polynomials at an integer index.
The bounds (I12) do not imply all-integer exclusion: no
Diophantine lower bound on $|m-\alpha_c|_5$, algebraicity or
non-Liouville property of $\alpha_c$ has been proved.
Valuations $v_5(u)\ne1$ in the already-confined inner region
$v_5(u)>3/5$ are outside this result.
No all-terminal quintic SFC, independent replay, external review,
formal verification, or novelty claim is made.


\section{A literal obstruction to forgetting the descendant seeds}
\begin{proposition}[Constant-unit seeds do not preserve mixed-pair noncollision]
\label{prop:quintic-seed-blind-obstruction}
The following stronger assertion is false: for every constant-unit
seed $h$ and every closed-class potential $F$, the generalized
ordinary quintic mixed pair has rank two at every nonzero argument
in the inner region.
In fact take $F=z$, $X=z-\theta_z$, $m=2$, and
\[
 \begin{split}
 d_0&=1349729058338,\\
 h(z)&=\frac{d_0+270777950790z-19101814440z^2
                         -542394950285z^3}{d_0}.
 \end{split}
\]
Then $h\in1+5\ZZ_5[z]$.  Define
\[
 B=(X+10)(X+7)h,\quad C=(X+13)B,\quad
 D=(X+16)C,\quad H=70D+3z^3B.
\]
At $u=z^3=10$, the two three-section rows of $(C,H)$ have
rank exactly one.
\end{proposition}
\begin{proof}
The denominator $d_0$ is a five-adic unit and each displayed
nonconstant numerator coefficient is divisible by five.
Direct integer polynomial multiplication gives
\[
 d_0(H-40C)=(z^3-10)R(z),
\]
where
\[
 \begin{aligned}
 R(z)={}&-132651371853458640\\
        &-62036179766657136z-11633463676444856z^2\\
        &-1088171938849410z^3-54296800471820z^4\\
        &-1627184850855z^5 .
 \end{aligned}
\]
Thus the two rows are proportional in
$\QQ_5[z]/(z^3-10)$.
The three sections of $d_0C$ evaluated at $u=10$ are
\[
 (-296770817069940,\ -112236517464324,\ -13462409209854),
\]
so the first row is nonzero and the rank is exactly one.
Since $v_5(10)=1>3/5$, this example lies in the inner region.
\end{proof}
This is a counterexample to a proposed stronger seed-blind invariant,
not an ordinary factorial-moment counterexample: its seed is
$h\ne1$.  The exact relationships among the descendant seeds in
the block transfer cannot be replaced by their constant-unit
reductions alone.

\section*{Verification scope}
The accompanying \href{quintic_five_adic_descent.py}{exact checker}
and \href{quintic_five_adic_descent_verification.json}{matching receipt}
check the complete finite offset transitions, exact polynomial
factorizations and symbolic residual identities, literal factorial
coefficients, seed and potential controls, rejection controls,
and Hensel correction digits.
The proofs above supply the stated unbounded quantifiers.
The finite controls alone do not prove those quantifiers.
No independent replay, external review, formal verification or
novelty assertion is made.
\end{document}
